\section{Numerical Methods and Evidence Standards}
\label{sec:numerical-methods}
\label{sec:numerical-methods-review}

Numerical evidence for stenotic-hemodynamics models is inseparable from the method that
produces the state, the reference data used to check it, and the operator used to
observe it. A resolved incompressible-flow or FSI computation, a distributed 1D
area--flow network, and a diagnostic postprocessing map can all report velocity
or flow, but they do so through different discrete spaces, stability
constraints, conservation properties, and measurement conventions. This chapter
therefore reviews numerical methods through the evidence standards needed to
interpret blood-flow simulations: verification, validation, stability,
positivity, conservation, well-balancing, reproducibility, and observation
operators.

The notation below separates the mathematical model from the discrete
realization:
\begin{flalign*}
    \quad
    \mathcal M(U)=0
    \quad\text{(continuous model)},\qquad
    M_h\dot U_h=R_h(U_h)
    \quad\text{(semidiscrete equations)},\qquad
    U_h^{n+1}=\Phi_{\Delta t,h}(U_h^n)
    \quad\text{(fully discrete scheme)}. &&
\end{flalign*}
A reported numerical quantity is then
\begin{flalign*}
    \quad
    Y_{h,\Delta t}=\mathcal O_h(U_h^n),
    &&
\end{flalign*}
with an observation map $\mathcal O_h$ that carries its own quadrature,
coordinate, and output conventions.

\subsection{Finite-Element Incompressible Flow and FSI}
\label{subsec:numerical-methods-fem-fsi}

Resolved hemodynamic models commonly start from incompressible Navier--Stokes or
Stokes equations, optionally coupled to a deforming wall. A finite-element
realization writes those equations in weak form, chooses discrete velocity and
pressure spaces, and solves the resulting saddle-point problem. The mixed
velocity--pressure pair is not a cosmetic choice: the discrete spaces must
control the pressure null mode and avoid spurious pressure oscillations. Stable
Taylor--Hood-type pairs, divergence-conforming constructions, or equal-order
pairs augmented by pressure-stabilization, streamline-upwind/Petrov--Galerkin,
least-squares, local projection, or grad-div terms are different numerical
contracts. Galdi and coauthors are useful here not as a generic FEM citation,
but because their hemodynamical-flow treatment ties continuum equations,
mixed discretization, stabilization, and simulation evidence together
\cite{GaldiEtAl2008HemodynamicalFlows}.

When the wall moves, the numerical contract also includes the fluid--structure
interface. FSI formulations must transfer traction and velocity or displacement
data across the interface, update the fluid domain or mesh, and control the
added-mass and time-lag effects that appear in partitioned coupling. Monolithic
and partitioned schemes can both be useful, but the FSI and multiscale coupling
sources show why they expose different solver errors, stability limits, and
coupling residuals
\cite{FormaggiaEtAl2007FSICoupling,QuarteroniVeneziani2003GeometricalMultiscale}.
For a reduced 1D comparison, a resolved FSI or CFD field is therefore not an
unqualified reference standard; it is a numerically produced state that still
needs its own mesh, timestep, boundary, material, and output-operator record.

\subsection{Finite Differences, Finite Volumes, DG, and High-Resolution Schemes}
\label{subsec:numerical-methods-balance-law-families}

One-dimensional blood-flow solvers are usually written as balance-law solvers.
Finite-difference schemes approximate derivatives at grid points and can be
compact and high-order, but conservation and geometric-source balance must be
checked at the discrete level. Finite-volume schemes instead advance cell
averages using numerical fluxes through cell faces. That structure makes local
conservation visible and is well suited to wave propagation, boundary fluxes,
and discontinuity-capturing limiters
\cite{LeVeque2002FiniteVolumeMethodsForHyperbolicProblems}. In arterial
networks, Wang's 1D network treatment adds the next layer: each vessel solve
must be compatible with node and terminal rules that determine incoming and
outgoing characteristic information \cite{Wang2014BloodFlowNetworks}. The
method comparison is therefore not finite volume versus network modeling; it is
how a conservative vessel update is embedded in a coupled vascular system.

Discontinuous Galerkin methods combine finite-volume fluxes with local
polynomial approximation. They can raise spatial order, expose element-local
residuals, and support $p$-refinement studies, but they still require numerical
fluxes, limiters or troubled-cell controls, and timestep choices. WENO and other
high-resolution reconstructions serve a related purpose for finite-volume and
finite-difference methods: they recover high-order accuracy on smooth regions
while suppressing nonphysical oscillations near steep gradients. In stenotic
area--flow problems, these design choices interact with wall and geometry
sources, so smooth-solution accuracy alone is not enough. A scheme can show good
manufactured-solution behavior and still fail to preserve a physiologically
meaningful rest state if flux and source terms are not balanced.

\subsection{Stability, Positivity, Conservation, and Well-Balancing}
\label{subsec:numerical-methods-stability-positivity-balance}

Explicit wave-propagation schemes carry CFL restrictions: the timestep must stay
compatible with the grid scale and characteristic speeds. Semi-implicit and
implicit network methods can relax the fastest explicit restriction, especially
when wave speed or network coupling is stiff, but they replace it with nonlinear
or linear solver tolerances, splitting errors, and coupling-residual criteria.
The benchmark and open-solver literature makes this evidentiary rather than
only algorithmic: Boileau and coauthors compare numerical schemes on shared 1D
arterial cases, OpenBF exposes a finite-volume solver record, and recent
well-balanced or Lax--Friedrichs studies show how boundary and coupling
policies affect the accepted run. Such workflows record not only the equation
family, but also the timestep policy, boundary assumptions, network coupling,
and diagnostic metadata
\cite{BoileauEtAl2015Benchmark,BenemeritoEtAl2024OpenBF,BeckersKolbe2025LaxFriedrichsHemodynamics,
LuccaEtAl2025FFRPrediction}.

For compliant-vessel balance laws, three checks are especially direct. First,
positivity guards ensure that area or radius states remain physical; any
projection or floor must be counted because it changes the discrete solution.
Second, conservation diagnostics record whether mass or flow bookkeeping is
consistent with boundary fluxes and source terms. Third, well-balanced tests
ask whether a target equilibrium is preserved by the discrete operator itself.
High-order well-balanced schemes make this requirement explicit by reconstructing
or integrating flux and source contributions so that selected rest states cancel
after discretization, not only in the continuous equations
\cite{PimentelGarciaEtAl2023HighOrderWellBalanced}. The 2023 JCP article
targets all stationary solutions of a model with discontinuous geometrical and
mechanical properties; the later friction-and-gravity preprint is a distinct
extension rather than the source for that claim. For the present report,
geometry-rest preservation is therefore a core numerical property rather than a
postprocessing preference or a cosmetic refinement of an otherwise verified
scheme.

\subsection{Verification, Validation, and Reference Data}
\label{subsec:numerical-methods-verification-validation}

ASME-style VVUQ terminology separates credibility claims
\cite{ASME2022VVUQTerminology}. Code verification concerns whether the
implementation conforms to the mathematical description. Solution verification
concerns numerical error in a computed quantity. Validation concerns whether a
model represents a specified real application for a specified observable and
context of use. Uncertainty quantification concerns the effect of uncertain
numerical and physical inputs. These categories are related, but none is a
substitute for another.

Verification asks whether the implementation solves the equations it claims to
solve. Manufactured solutions are useful because the analyst chooses a smooth
state, derives the forcing needed to make that state exact, and then checks the
discrete error under grid and timestep refinement. Exact or manufactured states
can probe fluxes, sources, time integration, boundary traces, and coordinate
maps without requiring physiological data. They also expose whether a nominally
high-order method is limited by boundaries, source quadrature, limiters, or
operator splitting.

Grid and timestep convergence studies then ask whether changes in resolution
change the reported quantities in the expected direction and magnitude. These
studies are strongest when spatial and temporal effects are separated; otherwise
they are better reported as sensitivity or time-step-insensitivity checks.
Benchmark problems provide a broader reproducibility context by comparing
multiple solvers on shared model cases. The Boileau benchmark is strongest in
that role: it helps interpret scheme behavior under common assumptions, but a
benchmark row does not by itself establish accuracy for a different stenosis
geometry, wall law, or boundary rule \cite{BoileauEtAl2015Benchmark}.

Open modeling platforms and model repositories occupy a related but distinct
evidence role. A platform such as SimVascular supports inspectable model
construction, meshing, simulation, and postprocessing workflows, while a
vascular model repository provides reusable geometry and model records
\cite{SimVascular2026Platform,SimVascular2026VascularModelRepository}. Those
resources strengthen reproducibility and make reference cases easier to inspect.
Their contribution is traceability of workflow and model records, not automatic
validation evidence; validation still requires the geometry, boundary data,
material assumptions, numerical settings, and observable to match the claim
being tested.

External validation is therefore a different evidence category. It requires
comparison with an accepted physical or application reference for the target
question, such as controlled experimental measurements or clinical measurements
with known acquisition conventions. A resolved computational comparator can
support a cross-model comparison when geometry, material parameters, boundary
history, timestep, and observable are matched closely enough for the intended
claim, but it is not automatically validation data merely because it is
higher-dimensional. In every category, the observation operator determines what
quantity is actually being compared.

\subsection{Cross-Dimensional Observation Operators}
\label{subsec:numerical-methods-observation-operators}

Reduced and resolved models do not naturally live in the same data space. A 1D
area--flow state supplies section area, flow, and mean axial velocity after a
coordinate convention is fixed. A 3D velocity field supplies nodal, elemental, or
quadrature-point velocities on a mesh. Comparing them therefore requires an
observation operator: a declared map from each numerical state to a common
observable.

For section comparisons, the operator must specify the target planes, mesh
intersection rule, interpolation of velocity onto the cut geometry, quadrature
weights, area normalization, and treatment of empty or invalid cuts. For profile
comparisons, it must also specify the radial coordinate, binning rule, profile
closure used by the 1D reconstruction, and any azimuthal averaging. Multiscale
and FSI literature treats such maps as part of the model interface rather than
as after-the-fact plotting choices
\cite{FormaggiaEtAl2007FSICoupling,QuarteroniVeneziani2003GeometricalMultiscale}.
Clinical FFR and CT-FFR sources show the same requirement in diagnostic form:
overview and CT-derived FFR work make the reported pressure-ratio observable
interpretable only after the imaging, model, boundary, and pressure-extraction
conventions have been fixed
\cite{TebaldiEtAl2015FFROverview,NorgaardEtAl2014NXTFFRCT,ShaikhEtAl2025CTFFR}.
Comparisons of 1D and 3D FFR estimates make the cross-dimensional version of
that requirement explicit \cite{Blanco2018FFR1D3D}.
In this report, the plane--tetrahedron quadrature operator plays that role for
the available resolved 3D velocity data, while the 1D state is observed through
physical flow and area-mean velocity.
A reported discrepancy can then include model-form error, closure or parameter
differences, numerical error, data-matching error, and observation-operator
error. The section geometry used by the observation map must also be declared:
reference-coordinate cuts and deformed-coordinate cuts are different output
operators even when they are applied to the same velocity file.

\subsection{Evidence Standard for the Present Report}
\label{subsec:numerical-methods-evidence-standards}

The review above gives a compact standard for reading the numerical record.
Each reported result should identify its equation tier, discrete operator,
boundary rule, timestep policy, stability and positivity diagnostics, reference
data, and observation operator. Table~\ref{tab:numerical-methods-evidence-map}
summarizes the evidence roles used by the present manuscript and later applied
in the case study.

\begin{table}[!htb]
    \centering
    \scriptsize
    \caption{Claim-to-evidence map for numerical-method statements in the
    report.}
    \label{tab:numerical-methods-evidence-map}
    \begin{tabularx}{\textwidth}{@{}p{0.42\textwidth}X@{}}
        \toprule
        \textbf{Claim} & \textbf{Minimum evidence} \\
        \midrule
        Implementation reproduces a manufactured smooth solution
            & Code verification and observed order for the declared forced
            operator, including flux, source, boundary, and time-integration
            wiring. \\
        Computed quantity is grid/time resolved
            & Solution verification with grid and timestep sensitivity or an
            uncertainty estimate for the reported observable. \\
        Rest or steady state is preserved
            & Named equilibrium family plus a drift, residual, or exact-discrete
            preservation test. \\
        Two schemes agree
            & Shared benchmark, matched model contracts, and comparable output
            quantities. \\
        Two model tiers agree
            & Matched inputs and closures plus a common observation operator. \\
        Model represents physical or clinical reality
            & External data for the specified observable and context of use. \\
        Result is reproducible
            & Archived or declared inputs, versions, commands, generated outputs,
            and source/provenance records. \\
        \bottomrule
    \end{tabularx}
\end{table}

\begin{stenosisillustration}
The later case study applies three of these evidence categories. The MMS run is
the manufactured-state verification problem: it checks whether the declared
finite-volume operator, forcing residuals, boundary traces, and timestepper
reproduce a known smooth area--flow solution. The rest-equilibrium run is the
well-balanced test: it asks whether the discrete stenosis geometry, wall
source, and numerical flux preserve the quiescent state that the continuous
balance law admits. The plane--tetrahedron quadrature rule is the
observation-operator test: it defines how a resolved 3D velocity field becomes
section area, physical flow, and area-mean velocity comparable to 1D outputs.
These checks answer different numerical questions, so agreement in one cannot
substitute for a limitation or ambiguity in another.
\end{stenosisillustration}
